\chapter[Discovery, Encounter and Conquest]{Were Ocean Voyages ``Discovery'', or Encounter, Conquest and Catastrophe?}
\section[A Spanish Dollar in China]{A Spanish Silver Dollar in a Chinese Cashbox}
Pick up an object, not from a museum display case, but from an old cashbox.

Imagine the Qianlong era of the Qing dynasty. At dusk, a tea merchant on the Fujian coast closes his stall and tips the day's takings into a wooden chest. Copper coins clatter, but several different ones are mixed in: much larger, silvery white and heavy, with a king's portrait on one side and a shield and castles on the other, their edges finely milled. These are ``foreign coins'', known locally by nicknames such as ``Buddha heads'' and ``foreign cakes''. Their proper name is benyang: Spanish silver dollars.

Here is the oddity. China did not make these coins, and the imperial court never approved them as legal tender, yet the whole south-eastern coast accepted them. Rice shops and cloth merchants took them; money shops kept special scales and silver-testing equipment for them. Their consistent fineness and full weight made them much easier to use than pieces of repeatedly clipped silver. Well into the nineteenth century, they remained among the coast's hardest currencies. An empire unable to mint such coins nevertheless filled its cashboxes with them.

Trace this coin further and its journey becomes astonishing. Its silver probably came from Potosí, high in present-day Bolivia: a ``silver mountain'' controlled by Spain from the mid-sixteenth century. Over the next two or three centuries, a substantial share of the world's newly mined silver came from that mountain. Shipped to Mexico City or Lima, it became large coins worth eight reales. A real was a Spanish monetary unit; ``eight reales'' was rather like saying ``a hundred-unit banknote''. The silver then took two routes. One went east across the Atlantic to Europe. The other went west on Spain's Manila galleons across the entire Pacific, unloading in the Philippines to purchase Chinese silk, porcelain and tea. Thus the coins entered Chinese cashboxes.

One coin connected four continents: American mines, European mints and royal portraits, Asian markets, and the whole Pacific linking them. This chain had not always existed. It had a starting point: October 1492, and three small ships from Spain.

Our textbooks call these voyages ``the opening of new sea routes'', formerly ``the great geographical discoveries''. The English-speaking world's formulation is most direct: Columbus ``discovered'' America. In this chapter, we shall weigh ``discovery''. The word can accommodate those three ships. Can it also accommodate the whole continent before their bows, and the tens of millions already living there?

\section{The Beach on 12 October 1492}
Come with me to a scene.

The Caribbean: a low island in what is now the Bahamas. It is the early morning of 12 October 1492. Three ships lie offshore. The largest, the Santa María, is a cumbersome cargo vessel; the Niña and Pinta are much smaller. Together they carry about ninety crew. They have seen no land for more than thirty days. The food is turning foul, the water growing algae, and grumbling verges on mutiny. The previous night brought drifting branches and birds not normally found far offshore. Around two in the morning, a sailor on the Pinta cries out: land.

Some of what happened after dawn we can establish; some we can only imagine. That distinction matters, and we shall return to it repeatedly.

What we can establish appears in the voyage journal. The original is lost, surviving through an abridgement copied decades later by a priest, Las Casas. It records islanders approaching the ships by canoe or swimming, almost unclothed, their bodies painted, some wearing small gold nose ornaments. They happily traded parrots, cotton thread and spears for glass beads and little bells. Columbus noticed that their spears were wooden shafts tipped with fish bones, that they did not know iron, and that they were generous. The journal's meaning is that whatever you asked for, they gave.

The other side requires imagination. Put yourself in an islander's body. You grew up fishing and growing cassava here; in your language the island is Guanahani. This morning three moving ``mountains'' appear on the horizon, covered in white cloth. Small boats descend from them, carrying people wrapped from head to foot in cloth and iron. You have never seen so much cloth, or any iron. They have pale skin, hairy faces and voices producing an entirely unintelligible stream of sounds. Your people offer food and water according to the rules of hospitality. You notice their leader staring at your gold nose ornaments, pointing and repeatedly saying something.

Two small events that day and in the following days are clearly recorded and worth remembering. The leader, the Genoese Columbus, declared the island the property of distant monarchs and renamed it San Salvador. He also took several islanders away, ostensibly to ``learn the language'' and serve as guides. Nobody asked for their consent, either to the renaming or to the removal of people.

We still know the name Guanahani only by good fortune: European records happened to mention it. The people they encountered became known as the Taíno. A few generations later, Taíno societies in the Caribbean had almost ceased to exist. That later development is the chapter's real issue. Before describing it, we must frame the question clearly.

\section{Do Not Rush to Call It ``Discovery''}
Let us set out the conflict without rushing to answer it.

On one side is a familiar account, some version of which appears in almost every student's history textbook. In the late fifteenth century, European navigators used courage, technology and curiosity to cross unknown oceans and connect two isolated continents: ``the world began to become a single whole''. Here 1492 marks a watershed, the beginning of modern history and an advance in human civilisation. ``Discovery'' seems entirely justified: European maps had not previously included America, which was indeed a ``New World'' for Europeans.

On the other side stands an objection so simple it almost seems sly: how could America be ``discovered'' when people already lived there? And not scattered groups of hunters, but an entire hemisphere of human societies.

To assess this objection's strength, consider its factual basis.

When Columbus landed, Tenochtitlan, the Aztec capital, stood in a lake on the site of today's Mexico City. Broad causeways, orderly waterways, a market serving tens of thousands daily and aqueducts carrying fresh water sustained it. Most researchers estimate its population in the hundred-thousands, up to around 200,000: larger than most contemporary European cities. Bernal Díaz, a Spanish soldier who first entered in 1519, later wrote in The True History of the Conquest of New Spain that they could scarcely believe their eyes and thought they were dreaming. Notice what astonished him: not ``backwardness'', but prosperity.

Farther south, another giant stretched across the Andes: the Inca Empire. It built roads along thousands of kilometres of mountain ridges and, without writing, managed taxation and inventories through quipu, knotted cords encoding information through colour, position and knot type. In Central America, the Maya had long possessed writing and precise calendars. In North America's Mississippi basin, Cahokia had built a city-state with enormous earth mounds and over ten thousand inhabitants centuries earlier. Estimates of the Americas' population before 1492 range widely, from under ten million to over a hundred million, but most accommodate a scale of tens of millions.

Here ``discovery'' meets its first hard obstacle: you can only discover something unknown to you. In ``Columbus discovered America'', Columbus is the subject and Europe the beneficiary; the tens of millions already there never appear. The sentence is not false---Europeans had not known America---but it is incomplete, and its incompleteness is not neutral.

Do not simply reverse the conclusion, either. If ``discovery'' is wrong, what is right? ``Encounter''? ``Invasion''? ``Conquest''? ``Catastrophe''? How much truth does each contain, and what problems does each conceal? ``Encounter'' sounds gentle and equal. But if one party loses most of its population while the other's monarchs grow rich, might that word be another whitewash? This is our real question: what should we call the events after 1492? Before answering, we must hear the different voices.

\section{Many Voices across the Compass}
\subsection[Why Europeans Went to Sea]{Why Europeans Went to Sea: Technology, Coins, Crowns and Crosses}
First hear Europe's reasons for itself: four intertwined reasons, not one.

\textbf{Technology.} None of Europe's navigational equipment was its invention alone. The caravel was a shallow-draught vessel capable of carrying triangular lateen sails from Arab and Indian Ocean traditions. These let ships tack in zigzags against the wind instead of waiting for a following breeze. The compass used a magnetic needle first employed for navigation in China, reaching the Mediterranean through Indian Ocean trade. The astrolabe and quadrant came through Islamic astronomy, estimating latitude from stars' altitude. The great voyages began as an assembly of Old World technologies.

\textbf{Capital.} Ocean voyages were fantastically expensive: ships, sailors, provisions and trade goods serving as ballast added up to a fortune. Where did it come from? Partly from Genoese and Florentine bankers and merchants pooling investments, sharing risks and dividing profits by shares. Private merchants held stakes in Columbus's voyage. From the outset, exploration was a business with written profit-sharing terms.

\textbf{Competition between states.} Under Prince Henry, remembered as ``the Navigator'', the Portuguese explored southwards along West Africa, seeking direct access to gold and spices without middlemen. Dias rounded the Cape of Good Hope in 1488; da Gama reached India in 1498. Spain had just captured Granada in January 1492, concluding the centuries-long Reconquista. Its treasury was depleted, while Portugal kept advancing into the Atlantic. The result was the Capitulations of Santa Fe in April 1492. The crown authorised Columbus's westward voyage, promising him the title ``Admiral of the Ocean Sea'' and a tenth of the proceeds if he found new lands. State rivalry transformed a Genoese sailor's private scheme into competition between crowns.

\textbf{Religion.} Spreading Christianity was an official objective and a sincere motive for many. The recently concluded Reconquista had embedded ``fighting for the faith'' in Spanish society's muscle memory. Distinguish two statements: believers sincerely considered missionary work good, and missionary work became a licence for conquest. Both can be true. Understanding a motive does not excuse its consequences.

Here is \textbf{a counterexample easy to misjudge}. You have probably heard the popular claim that the Ottomans' capture of Constantinople in 1453 cut East--West trade routes, forcing Europeans to seek alternatives by sea. It makes a smooth story and appears in countless popular books, but does not withstand scrutiny. Trade was not cut off: spices continued through Egypt and the Red Sea to Venetian merchants. The Portuguese motive was less dramatic and more powerful: numerous middlemen multiplied spice prices before they reached Europe, and bypassing them promised riches. The ``blocked trade routes'' story resembles a dramatic prologue supplied afterwards. Remember: an explanation may become popular simply because it tells well.

\subsection[The Strongest Case for Commemoration]{Making the ``Discoverers' '' Case in Full}
Before criticising further, practise making the strongest case for commemorating Columbus: one his supporters would themselves endorse.

Here it is. The courage was real: he sailed into waters maps dared not depict and held a nearly mutinous crew to a life-or-death gamble. The skill was real: using stars, winds and currents, he made four return voyages. His trade-wind route remained standard for centuries. Most importantly, the connection was real. After 1492, hemispheres separated for over ten thousand years began exchanging things. Maize, potatoes, sweet potatoes, tomatoes, chillies and cocoa spread from the Americas worldwide; high-yielding crops subsequently fed countless people in Eurasia and Africa. Horses, cattle, wheat and sugar cane travelled to America. Historian Alfred Crosby called this the ``Columbian Exchange''. From that perspective, 1492 was a genuine watershed for all humanity. No islanders' catastrophe can negate the lives later saved by potatoes and sweet potatoes in the Old World.

This is commemoration's strongest version: honouring the connection, not Columbus's character. His abuse and enslavement of islanders became intolerable even to the contemporary Spanish crown, which sent him home in chains in 1500.

Remember the weight of this argument. The voices that follow need not erase it; they ask who is absent from its narrative.

\subsection[The Taíno]{The Taíno: First to Welcome Them, First to Disappear}
Return to the Caribbean. Within decades of 1492, Taíno society on Hispaniola, today's Haiti and Dominican Republic, collapsed. We shall examine how later. First hear a voice that corrects a common picture: the Taíno were not passive scenery waiting to be attacked.

In 1511, the Taíno leader Agüeybaná II led resistance in Puerto Rico. Later chronicles report an earlier ``experiment'': the Taíno held a Spaniard underwater in a river until they established that these supposedly supernatural foreigners could drown. Then they acted. Retellers may have embellished the details, but the story preserves something important: Taíno agency. They observed, judged, tested and chose to fight.

The resistance was suppressed. Survivors fled to the mountains, later joining escaped African slaves. Many Caribbean people today have Taíno ancestry, while Taíno words survive worldwide: ``hurricane'', ``canoe'', ``hammock'' and ``barbecue'' passed through Spanish into English and beyond. A people described as having ``disappeared'' is living in your dictionary.

\subsection[Conquest in the Valley of Mexico]{The Valley of Mexico: More Than a Few Hundred Conquerors}
In 1519, Cortés landed in Mexico with a few hundred soldiers. Two years later, Tenochtitlan fell. For centuries, the standard question has been: how did a few hundred Spaniards defeat an empire?

The question's wording contains a clue. The answer is: they did not do it alone. Cortés quickly discovered that the Aztec Empire was a tribute-based hegemony deeply resented by neighbouring city-states, especially the Tlaxcalans. Spaniards supplied steel weapons, horses and an anti-Aztec banner. Tlaxcalans supplied tens of thousands of warriors, food, porters and complete knowledge of the terrain. Tenochtitlan fell to a predominantly Indigenous coalition. The Spaniards were its smallest but least replaceable component.

Clarifying this does not excuse the Spaniards: they commanded, betrayed and massacred after the city's fall. It corrects a myth useful to later propaganda on both sides. Spaniards celebrated victory against overwhelming odds as evidence of divine aid. The tragic version, ``a few hundred destroyed a country'', likewise erases Tlaxcalan choices. They fought for their own interests, winning their wager on one battle but misjudging the world afterwards. As meritorious allies, they received privileges. Within decades, however, they found their new masters far greedier than the old.

\subsection{Earlier Ocean Crossers}
Finally hear three often omitted voices. They show that 1492 was neither humanity's first ocean crossing nor the only way to cross an ocean.

\textbf{The Norse.} Around 1000 CE, Norse voyagers from Greenland and Iceland reached Newfoundland in present-day Canada and settled at L'Anse aux Meadows. Archaeologists excavated the unquestionable remains in the 1960s. In 2021, scientists used a solar storm's tree-ring signature to date the felling of several trees to 1021. Columbus was not the first European in America. But that encounter produced no sustained exchange; the small settlement was abandoned after some years. Between arrival and changing the world lie population, capital, states and supply lines. Columbus's ``first'' was not arriving, but being caught by an entire supporting system that never let go.

\textbf{The Polynesians.} The Pacific provides an even stronger comparison. Without compasses or charts, Polynesian navigators used rising and setting stars, currents, cloud formations and seabird routes to sail double-hulled canoes throughout the triangle bounded by New Zealand, Hawai‘i and Easter Island. Thousands of kilometres of empty ocean separated landfalls. Marshallese navigators made ``stick charts'' from palm ribs and shells, encoding reflected and interfering waves between islands in portable teaching aids. In 1976, Hawaiians built the traditional canoe Hōkūle‘a. The veteran Micronesian navigator Mau Piailug guided it from Hawai‘i to Tahiti, thousands of kilometres away, without modern instruments. It arrived, demonstrating the reality of that knowledge. Remember this vessel if you imagine ocean voyaging a European monopoly.

\textbf{Zheng He.} Between 1405 and 1433, decades before Columbus's birth, Ming fleets made seven western voyages. Over a hundred ships and more than twenty thousand people made Columbus's three vessels look like a handful of dinghies. They reached East Africa. Then the voyages stopped. Historians debate why: excessive cost, unprofitable tribute trade, or a shift towards northern frontier defence? The contrast is striking. Zheng He's fleets did not seize ports or enclose territory; they exchanged gifts and established tributary relations. European vessels after 1492 carried contracts, cannon and profit-sharing clauses. This often supports the claim that ``Chinese people were peaceful, Europeans greedy'', but that is too slippery. Tribute was also power; Zheng He's fleet fought in Sri Lanka and captured a king. More honestly, the expeditions sought different things: prestige and order in one case, profit and territory in the other. The latter carried an engine-like logic of expansion which, once started, would not stop by itself.

\section[Thinking Bridge: Change the Subject]{Thinking Bridge: Change the Subject and Retell the History}
After these voices, we need a portable tool. Its central insight is this: \textbf{a historical narrative's viewpoint hides in its grammar, specifically its subject.}

There are four steps:

\begin{enumerate}
\item Find the sentence's subject: who is acting?
\item Ask who else participates in or experiences the action. List everyone, omitting nobody.
\item Make someone else the subject and rebuild the sentence, changing the verb if necessary.
\item Compare versions. Which facts remain true in all? Which become visible only in one?
\end{enumerate}
Try it with our main example:

\begin{itemize}
\item \textbf{Version one:} ``In 1492, Columbus discovered America.'' Columbus is the subject. Check: true within Europe's knowledge system, but tens of millions vanish from the sentence.
\item \textbf{Version two:} ``In 1492, the Taíno encountered uninvited strangers carrying iron and contracts.'' The Taíno are the subject. Equally true, it reveals something else: subsequent events become ``responding to invasion'', not ``welcoming progress''.
\item \textbf{Version three:} ``After 1492, two hemispheres isolated for over ten thousand years began exchanging species, germs and silver.'' Exchange is the subject. Both sides become participants, accommodating the Columbian Exchange. But beware its excessive gentleness: ``exchange'' conceals whose guns determined the terms.
\item \textbf{Version four:} ``In 1520, smallpox conquered Tenochtitlan; the Spanish coalition followed a year later.'' The virus is the subject. Strange-sounding but informative, this version will matter in the next section.
\end{itemize}
Practise on any historical statement. In ``Zheng He sailed the Western Seas'', the subject could be Zheng He, the Yongle emperor, a Fujian sailor or a merchant in Malacca. Each sees a different voyage. Try a contemporary news headline with equally striking results.

Set a boundary, however: changing the subject reveals \textbf{perspective}, not a substitute for \textbf{evidence}. ``The Taíno met strangers'' and ``Columbus discovered America'' are grammatically equal, but not equal in historical weight: one account's consequences included a hemisphere-wide population catastrophe. Changing subjects reveals the parties; afterwards you must weigh them, using evidence. Let us read some.

\section{A Letter Sent Home from the Sea}
On his return voyage in 1493, Columbus wrote to the royal financial official Santángel, reporting his first expedition. Printed in Barcelona that year and translated into Latin within months, the letter circulated through European courts in more than a dozen editions: an early viral news story. We can still read it.

He reported many things: discovering numerous islands and taking possession of each through ceremonies involving royal flags and proclamations. One sentence deserves particular attention. When occupying these islands, he said:

\begin{quote}
``Nobody objected.'' (Paraphrased.)
\end{quote}
In the journal's October 1492 entries, preserved through Las Casas's abridgement, Columbus judged that the islanders should make ``very good servants'' and could ``easily become Christians''. At the letter's end he promised the crown as much gold, spice and enslaved labour as it requested.

Apply our tool. Notice three adjacent verbs: \textbf{occupy, name, promise}. On the beach, these acts occurred almost together. Consider the extraordinary ``Nobody objected''. Naturally nobody did: an objection requires shared language and a shared court, neither available to the islanders. A judgement was delivered in a courtroom where only one party was present, then recorded as unopposed. ``Good servants'' and ``Christians'' sit comfortably together in Columbus's thinking: making people servants and offering their souls to God were both ``for their good''. Understanding does not mean accepting this. But it helps explain what followed. The Caribbean catastrophe did not begin when someone suddenly became wicked; the letter's logic was amplified and implemented over decades.

A contrasting source adds weight. Decades later, Las Casas, a Spanish priest who spent half his life in America, wrote A Short Account of the Destruction of the Indies, published in 1552, reporting to the king what he had witnessed. He accused his countrymen of turning island after island into places of death through mining and forced labour. Most historians today consider his tolls of millions or tens of millions exaggerated. But his identity prevents dismissal as mere ``enemy propaganda'': he was the conquerors' compatriot, the king's cleric, an insider. A letter and a book from two members of one civilisation: one made people resources, the other victims crying for justice. Sometimes the most painful discovery in hearing multiple voices is that disagreement divides a civilisation internally, not merely one civilisation from another.

\section[Where Did the Millions Go?]{Where Did Tens of Millions Go? Three Explanations}
Now we face the chapter's heaviest page, where distinguishing our four kinds of content matters most.

\textbf{What happened?} At the evidential level, little is disputed: Indigenous American populations collapsed after contact. The severity varied, but many regions lost over half their people within a century. The Caribbean suffered worse, with entire societies disintegrating in a few generations. A hemisphere-wide total is impossible: scholars cannot even agree on the starting population. But there is no scholarly dispute over the scale: one of human history's greatest demographic catastrophes.

\textbf{Why did it happen?} At the explanatory level, three accounts contend. Consider each one's arguments and limitations.

\textbf{Explanation one: germs, or ecological epidemiology.} Isolated from the Old World for over ten thousand years, Indigenous Americans lacked immunity to smallpox, measles and influenza, long ``domesticated'' into childhood diseases in Eurasia and Africa. Smallpox entered Mexico with Spanish forces in 1520 and outran them. It killed far more people than firearms, including leaders, priests and record-keepers---society's brain. The Inca emperor Huayna Capac very probably died of imported smallpox; the ensuing succession war opened a door for Pizarro. This explains the collapse's speed and extent, including regions no European army had reached: germs travelled faster. Its limitation is equally clear. Alone, it can become a ``natural disaster'' narrative, a plague for which nobody need apologise. Yet decisions by people shaped how pathogens spread, why patients lacked care and why surviving communities were sent into mines rather than allowed to recover.

\textbf{Explanation two: violence, in the tradition of Las Casas.} War, massacre, enslavement and forced labour directly killed people. One institution needs explanation: the encomienda. The crown ``entrusted'' a territory's Indigenous population to colonists, ostensibly for protection and religious instruction. In practice this often legalised extracted labour. Abundant contemporary testimony supports the account; conquistadors' own letters and books record slave prices and labour quotas. It restores responsibility and fits eyewitness evidence. But blades, guns and whips cannot explain the whole scale: death reached places Spanish power would enter only decades later, or never rule directly.

\textbf{Explanation three: systemic collapse, the synthesis of most contemporary scholars.} No single cause sufficed. Germs, warfare, enslavement, famine and social disintegration interlocked like a mincing machine. Epidemics killed farmers, causing hunger; refugees carried infection farther; colonists exploited weakness to demand labour; malnutrition increased vulnerability to the next epidemic. This best fits the overall evidence and regional differences. Yet it carries a hidden risk: ``systemic'' can wash everything grey. The machine is a system, but someone switches it on, someone feeds it and someone takes the profits. Explaining systems must not erase responsible people.

As for our fourth category, \textbf{how we judge}, wait. You now have equipment: the fact of collapse; three explanations with strengths and weaknesses; and contemporary understandings. Columbus made people resources, Las Casas made them victims, and the Tlaxcalans saw a wager gone wrong. The judgement is yours at the chapter's end.

\section[Further Challenge: Maps as Claims]{Further Challenge: Is a Map a Mirror or a Claim?}
One question remains: how do maps and naming take possession of space? We need critical cartography, the study of how maps speak.

Consider three things. First, in 1493 Pope Alexander VI issued a bull drawing a north--south Atlantic line. Lands west of it not already held by Christian monarchs were ``given'' to Spain; those east went to Portugal. The 1494 Treaty of Tordesillas moved the line several hundred leagues west. With a pen, two countries divided an unseen world of unknown size across an ocean. That line explains why Brazil speaks Portuguese and most of the rest of Latin America Spanish.

Second, naming. Columbus renamed each island he reached: Guanahani became San Salvador. In 1507, German cartographer Waldseemüller adopted the account in navigator Amerigo Vespucci's letters and labelled the new continent ``America'', a Latin feminine form of Amerigo. Once printed on a map, a name could not be peeled away. An entire hemisphere now bears a Florentine's name, while the Taíno Guanahani survives in footnotes.

Third, theory. Late-twentieth-century cartographers, notably Brian Harley, argued that \textbf{a map is never a neutral mirror; it is a claim}. Every map states what exists, who owns what, what matters and what does not exist. Its strongest weapon is often not a drawn feature, but blank space. Colonial maps regularly rendered inhabited lands empty. In contemporary European legal thinking, terra nullius---``nobody's land''---could legitimately be occupied. Erasing people was conquest's first step, accomplished in ink.

Notice the other side: erased people also made maps. We have encountered Marshallese stick charts. Indigenous hunters on Greenland's eastern Arctic coast carved coastal outlines in wood, carrying them inside their clothing and reading by touch in polar darkness. Today, Amazonian communities use satellite positioning to map their territories and submit maps in land cases, turning the colonisers' weapon around. If maps make claims, anyone can make one.

This theory needs a boundary too. Critical cartography goes too far if it says, ``Maps are nothing but power; truth and falsehood do not apply.'' Coastlines can be right or wrong. A chart thirty degrees out in longitude can wreck a ship. More precisely, reality constrains the geographical skeleton, but names, boundaries and blank spaces involve choices, and choices involve interests. Learn to see both skeleton and choices.

\section{The Encounter Is Not Over}
Echoes of this old question are in the news now.

\textbf{A tug-of-war on the calendar.} In the United States, Columbus Day became a national commemoration at the 1892 four-hundredth anniversary, largely symbolising recognition sought by discriminated-against Italian immigrants: Columbus was Genoese. Since the five-hundredth anniversary in 1992, growing numbers of states and cities have called the date Indigenous Peoples' Day. From 2021, presidents also proclaimed Indigenous Peoples' Day federally. Two commemorations pull at the same calendar square. Resistance to renaming comes not only from nostalgia for conquest, but from hurt Italian Americans whose excluded ancestors treated Columbus as proof that ``we belong here too''. Changing a holiday touches two communities' memories. Practise our skill: hear both arguments fully before judging.

\textbf{Statues and place names.} In 2020, Columbus statues in several US cities were toppled or officially removed. Mexico City's statue on Paseo de la Reforma was removed, with replacement plans debated for years. In 2015, North America's highest peak officially changed from Mount McKinley back to the Indigenous Alaskan name Denali, meaning ``the great mountain''. Renaming is no recent invention, and naming is never final. Names keep moving; 1492 was one particularly violent diversion.

\textbf{Textbooks change their language.} Chinese textbooks now generally prefer ``opening new sea routes'', already more restrained than ``great geographical discoveries''. European and American books increasingly use frameworks such as an ``age of contact'' and the ``Columbian Exchange'', placing both parties in one paragraph. Words change slowly, but changing them changes the telling of historical facts. Our thinking bridge is already operating in textbooks.

\textbf{Map battles have moved online.} Korean maps label a sea the ``East Sea'', Japanese maps the ``Sea of Japan''; the International Hydrographic Organization has debated it for decades without resolution. Navigation software shows different street names in different countries. Five centuries later, ``maps make claims'' remains news.

Another echo sits on supermarket shelves. Your potatoes, maize, chillies, tomatoes and chocolate came from the other hemisphere through that exchange. The encounter's consequences inhabit your body as well as monuments. This does not soften history, but makes ``us and them'' harder to say. Five hundred years later, almost everyone is a child of that exchange, at least at the dinner table.

\section[Your Turn: Whose Name for an Island?]{Your Turn: Whose Name Should an Island Bear?}
Return to the island where this chapter began. Its official name remains San Salvador, the name Columbus gave it over five centuries ago.

Imagine you are a fifteen-year-old resident, and the island holds a referendum on renaming. Three options appear: keep San Salvador, restore Guanahani, or choose an entirely new name. Which gets your vote?

Answer with reasons strong enough to address at least one objection:

Can renaming change what happened? Does keeping the name continue signing that morning's declaration? Can a name settle history's debts? If not, what is commemoration for?

Think a step further for your opponent. If you choose Guanahani, a classmate descended from Spanish sailors says, ``My grandfather called it San Salvador too. It was his home as well.'' What do you say? If you keep the name, a classmate identifying as a Taíno descendant says, ``Your name is a receipt for occupation.'' How do you answer?

There is no standard answer. Keep two things: a tool---find the subject of any historical statement---and an alertness---every name on a map was once a claim. Next time you look up at a world map, try to see the other names pressed beneath those printed there.
